% Part: incompleteness
% Chapter: representability-in-q
% Section: c-representable

\documentclass[../../../include/open-logic-section]{subfiles}

\begin{document}

\olfileid[pt]{inc}{req}{cre}
\olsection{Funções em $C$ são Representáveis em $\Th{Q}$}

Precisamos mostrar que toda função em $C$ é representável em $\Th{Q}$.
Ao final, precisamos mostrar como atribuir a cada função $k$-ária
$f(x_0,\dots,x_{k-1})$ em $C$ uma fórmula
$!A_f(x_0,\dots,x_{k-1},y)$ que a represente.

\begin{lem}
Dados números naturais $n$ e $m$, se $n \neq m$, então $Q \Proves \num
n \neq \num m$.
\end{lem}

\begin{proof}
Use indução sobre $n$ para mostrar que, para todo $m$, se $n \neq m$, então
$Q \Proves \eq/[\num n][\num m]$.

No caso base, $n = 0$. Se $m$ não é igual a $0$, então $m = k + 1$ para
algum número natural $k$. Temos um axioma que afirma
$\lforall[x][\eq/[0][x']]$. Por um axioma dos quantificadores, substituindo
$x$ por $\num k$, podemos concluir $\eq/[0][\num k']$. Mas $\num k'$ é
simplesmente $\num m$.

No passo indutivo, podemos supor que a afirmação seja verdadeira para $n$ e
considerar $n+1$. Seja $m$ um número natural qualquer. Há duas
possibilidades: ou $m = 0$, ou, para algum $k$, temos $m = k+1$. O primeiro
caso é tratado como acima. No segundo caso, suponha $n+1 \neq k+1$. Então
$n \neq k$. Pela hipótese de indução para $n$, temos
$\Th{Q} \Proves \eq/[\num n][\num k]$. Temos um axioma que afirma
$\lforall[x][\lforall[y][\eq[x'][y'] \lif \eq[x][y]]]$. Usando um axioma
dos quantificadores, temos $\eq[\num n'][\num k'] \lif \eq[\num n][\num
  k]$. Pela lógica proposicional, podemos concluir, em $\Th{Q}$,
$\eq/[\num n][\num k] \lif \eq/[\num n'][\num k']$. Usando modus ponens,
podemos concluir $\eq/[\num n'][\num k']$, que é o que queremos, pois
$\num k'$ é $\num m$.
\end{proof}

Observe que o lema não diz muita coisa: em essência, afirma que $\Th{Q}$
pode demonstrar que numerais diferentes denotam objetos diferentes. Por
exemplo, $\Th{Q}$ demonstra $0'' \neq 0'''$. Mas mostrar que isso vale em
geral exige algum cuidado. Observe também que, embora estejamos usando
indução, trata-se de indução \emph{fora} de $\Th{Q}$.

Seremos capazes de representar o zero, o sucessor, a soma, o produto, a
função característica da igualdade e as projeções. Em cada caso, a função
representante apropriada é inteiramente direta; por exemplo, o zero é
representado pela fórmula
\[
y = 0,
\]
o sucessor é representado pela fórmula
\[
x_0' = y,
\]
e a soma é representada pela fórmula
\[
x_0 + x_1 = y.
\]
O trabalho consiste em mostrar que $\Th{Q}$ pode demonstrar as afirmações
relevantes; por exemplo, dizer que a soma é representada pela fórmula acima
envolve mostrar que, para todo par de números naturais $m$ e $n$, $\Th{Q}$
demonstra
\[
\eq[\num n + \num m][\num {n+m}]
\]
e
\[
\lforall[y][(\eq[(\num n + \num m)][y] \lif \eq[y][\num {n+m}])].
\]

E quanto à composição? Suponha que $h$ seja definida por
\[
h(x_0,\dots,x_{l-1}) = f(g_0(x_0,\dots,x_{l-1}), \dots,
g_{k-1}(x_0,\dots,x_{l-1})).
\]
onde já encontramos fórmulas $!A_f, !A_{g_0}, \dots,
!A_{g_{k-1}}$ que representam, respectivamente, as funções
$f,g_0,\dots,g_{k-1}$. Podemos então definir uma fórmula $!A_h$ que
represente $h$, definindo $!A_h(x_0,\dots,x_{l-1},y)$ como
\begin{multline*}
  \lexists[z_0\dots][\lexists[z_{k-1}][(!A_{g_0}(x_0,\dots,x_{l-1},z_0) \land
  \dots \land !A_{g_{k-1}}(x_0,\dots,x_{l-1},z_{k-1}) \land]]\\
  !A_f(z_0,\dots,z_{k-1},y)).
\end{multline*}

Por fim, consideremos a busca ilimitada. Suponha que $g(x,\vec z)$ seja
regular e representável em $\Th{Q}$, digamos pela fórmula $!A_g(x,\vec
z,y)$. Seja $f$ definida por $f(\vec z) = \mu x \; g(x,\vec z)$. Gostaríamos
de encontrar uma fórmula $!A_f(\vec z,y)$ que represente $f$. Eis uma
escolha natural:
\[
!A_f(\vec z,y) \ident !A_g(y,\vec z,0) \land \lforall[w][(w < y
\lif \lnot !A_g(w,\vec z,0))].
\]

Pode-se mostrar que isso funciona usando alguns lemas adicionais, por
exemplo:

\begin{lem}
  Para toda variável $x$ e todo número natural $n$, $\Th{Q}$ demonstra $\eq[(x'
  + \num n)][(x + \num n)']$.
\end{lem}

Vale novamente mencionar que isso é mais fraco do que dizer que $\Th{Q}$
demonstra $\lforall[x][\lforall[y][(x' + y) = (x +y)']]$ (o que é falso).

\begin{proof}
A demonstração é, como de costume, por indução sobre $n$. No caso base,
$n = 0$, precisamos mostrar que $\Th{Q}$ demonstra
$\eq[(x'+0)][(x+0)']$. Mas temos:
\begin{eqnarray*}
& & \eq[(x' + 0)][x'] \quad \text{pelo axioma 4} \\
& & \eq[(x + 0)][x] \quad \text{pelo axioma 4} \\
& & \eq[(x+0)'][x'] \quad \text{pela linha 2} \\
& & \eq[(x' + 0)][(x + 0)'] \quad \text{linhas 1 e 3}
\end{eqnarray*}
No passo indutivo, podemos supor que derivamos $\eq[(x' +
  \num n)][(x + \num n)']$ em $\Th{Q}$. Como $\num{n+1}$ é $\num n'$,
precisamos mostrar que $\Th{Q}$ demonstra $\eq[(x' + \num n')][(x + \num
n')']$. A seguinte cadeia de igualdades pode ser derivada em $\Th{Q}$:
\begin{eqnarray*}
(x' + \num n') & \eq & (x' + \num n)' \quad \text{axioma 5} \\
& \eq & (x + \num n')' \quad \text{pela hipótese de indução}
\end{eqnarray*}
\end{proof}

\begin{lem}
\begin{enumerate}
\item $\Th{Q}$ demonstra $\lnot (x < \num 0)$.
\item Para todo número natural $n$, $\Th{Q}$ demonstra
\[
x < \num {n+1} \lif (\eq[x][\num 0] \lor \dots \lor \eq[x][\num n]).
\]
\end{enumerate}
\end{lem}

\begin{proof}
Façamos 1 e parte de 2 informalmente (isto é, dando apenas indicações de
como construir a !!{derivation} formal).

Para a parte 1, pela definição de $<$, precisamos demonstrar $\lnot
\lexists[y][\eq[(y'+x)][0]]$ em $\Th{Q}$, o que é equivalente (usando os
axiomas e as regras da lógica de primeira ordem) a
$\lforall[y][\eq/[(y' + x)][0]]$. Eis a ideia: suponha
$\eq[(y'+x)][0]$. Se $x$ é 0, temos $\eq[(y' + 0)][0]$. Mas, pelo axioma 4
de $\Th{Q}$, temos $\eq[(y' + 0)][y']$, e, pelo axioma 2, temos
$\eq/[y'][0]$, uma contradição. Logo, $\lforall[y][\eq/[(y' +x)][0]]$.
Se $x$ não é $0$, pelo axioma 3 existe um $z$ tal que $\eq[x][z']$. Mas
então temos $\eq[(y' + z')][0]$. Pelo axioma~5, temos
$\eq[(y' + z)'][0]$, novamente contradizendo o axioma~2.

Para a parte 2, use indução sobre $n$. Consideremos o caso base, quando
$n =0$. Nesse caso, precisamos mostrar $x < \num 1 \lif \eq[x][\num
  0]$. Suponha $x < \num 1$. Então, pelo axioma que define $<$, temos
$\lexists[y][\eq[(y'+x)][0']]$. Suponha que $y$ tenha essa propriedade;
assim, temos $\eq[y'+x][0']$.

Precisamos mostrar $\eq[x][0]$. Pelo axioma 3, se $x$ não é 0, ele é igual
a $z'$ para algum $z$. Então temos $\eq[(y' + z')][0']$. Pelo axioma~5 de
  $\Th{Q}$, temos $\eq[(y'+z)'][0']$. Pelo axioma 1, temos $\eq[(y' +
    z)][0]$. Mas isso significa, por definição, que $z < 0$, contradizendo
  a parte~1.
\end{proof}

\begin{explain}
Mostramos que o conjunto das funções computáveis pode ser caracterizado
como o conjunto das funções representáveis em $\Th{Q}$. Na verdade, a
demonstração é mais geral. Pela definição de representabilidade, não é
difícil ver que qualquer teoria que estenda $\Th{Q}$ (ou na qual se possa
interpretar $\Th{Q}$) pode representar as funções computáveis; mas,
reciprocamente, em qualquer sistema de !!{derivation} no qual a noção de
demonstração seja computável, toda função representável é computável.
Assim, por exemplo, o conjunto das funções computáveis pode ser
caracterizado como o conjunto das funções representadas na aritmética de
Peano ou até mesmo na teoria dos conjuntos de Zermelo--Fraenkel. Como
G\"odel observou, isso é um tanto surpreendente. Veremos que, no que diz
respeito à demonstrabilidade, as questões são muito sensíveis à teoria
considerada; grosso modo, quanto mais fortes os axiomas, mais se pode
demonstrar. Contudo, em uma ampla variedade de teorias axiomáticas, as
funções representáveis são exatamente as computáveis.
\end{explain}

\end{document}
